\chapter{طیف سیاسی در جهان غیرغربی}

\begin{tcolorbox}[
    enhanced,
    colback=lightbg,
    colframe=goldaccent,
    boxrule=2pt,
    arc=8pt,
    fonttitle=\bfseries\large,
    title={\rl{نقل‌قول آغازین}},
    attach boxed title to top center={yshift=-3mm},
    boxed title style={colback=goldaccent, colframe=goldaccent}
]
\begin{center}
\Large\textit{«ما نه شرقی هستیم نه غربی؛ ما مصری هستیم.»}

\vspace{3mm}
\normalsize
--- جمال عبدالناصر، ۱۹۵۶
\end{center}
\end{tcolorbox}

\vspace{5mm}

\section*{درآمد}

\begin{multicols}{2}
تا اینجا عمدتاً از تاریخ و اندیشهٔ سیاسی غرب سخن گفته‌ایم. این طبیعی است: مفاهیم «چپ» و «راست» در انقلاب فرانسه زاده شدند و در بستر تاریخ اروپا و آمریکا توسعه یافتند. اما آیا این مفاهیم در جهان غیرغربی هم کاربرد دارند؟

در این فصل به بررسی طیف سیاسی در مناطق مختلف جهان می‌پردازیم: خاورمیانه و جهان عرب، ایران، آسیای جنوبی و شرقی، آمریکای لاتین، و آفریقا. خواهیم دید که گاه مفاهیم غربی قابل انطباق‌اند، گاه نیازمند تعدیل، و گاه کاملاً نامناسب.
\end{multicols}

%═══════════════════════════════════════════════════════════════
\section{آیا چپ و راست جهانشمول است؟}
%═══════════════════════════════════════════════════════════════

\begin{tcolorbox}[
    enhanced,
    colback=leftred!8,
    colframe=leftred,
    boxrule=1.5pt,
    arc=5pt,
    fonttitle=\bfseries,
    title={\rl{🔴 مناظره: جهانشمولی یا غرب‌محوری؟}}
]
\begin{multicols}{2}
\textbf{چپ و راست جهانشمول است:}
\begin{itemize}
    \item تقسیم‌بندی بر اساس برابری جهانی است
    \item همهٔ جوامع نابرابری دارند
    \item مفاهیم قابل ترجمه‌اند
    \item جنبش‌های چپ و راست در همه‌جا وجود دارند
\end{itemize}

\columnbreak

\textbf{چپ و راست غرب‌محور است:}
\begin{itemize}
    \item تاریخ خاص اروپایی دارد
    \item شکاف‌های دیگر مهم‌ترند (دین، قوم)
    \item تحمیل چارچوب غربی
    \item نیاز به مفاهیم بومی داریم
\end{itemize}
\end{multicols}

\vspace{3mm}
\textbf{موضع این فصل:} مفاهیم چپ و راست تا حدی قابل تعمیم‌اند، اما باید با شکاف‌های محلی ترکیب شوند. در هر منطقه، شکاف‌های خاصی (دین، قوم، استعمار) اهمیت دارند که در چارچوب غربی نمی‌گنجند.
\end{tcolorbox}

%═══════════════════════════════════════════════════════════════
\section{خاورمیانه و جهان عرب}
%═══════════════════════════════════════════════════════════════

\subsection{ناسیونالیسم عربی}

\begin{multicols}{2}
در قرن بیستم، \textbf{ناسیونالیسم عربی} (پان‌عربیسم) جنبش غالب در خاورمیانه بود. این جنبش می‌خواست ملت‌های عرب را متحد کند.

\textbf{ناصریسم:}
\textbf{جمال عبدالناصر} (۱۹۱۸-۱۹۷۰) رهبر مصر از ۱۹۵۶ تا ۱۹۷۰ بود. او:
\begin{itemize}
    \item کانال سوئز را ملی کرد
    \item با غرب و اسرائیل مقابله کرد
    \item اصلاحات ارضی انجام داد
    \item صنایع را دولتی کرد
    \item با سوسیالیسم عربی تعریف شد
\end{itemize}

\textbf{بعثیسم:}
حزب بعث («رستاخیز») در ۱۹۴۷ تأسیس شد با شعار «وحدت، آزادی، سوسیالیسم». این حزب در عراق (صدام حسین) و سوریه (حافظ اسد) به قدرت رسید.

این جنبش‌ها چپ بودند یا راست؟ از نظر اقتصادی (دولت‌گرایی، اصلاحات ارضی) چپ بودند؛ اما از نظر اجتماعی محافظه‌کار و از نظر سیاسی اقتدارگرا.
\end{multicols}

\subsection{اسلام‌گرایی: کجای طیف؟}

\begin{tcolorbox}[
    enhanced,
    colback=leftred!8,
    colframe=leftred,
    boxrule=1.5pt,
    arc=5pt,
    fonttitle=\bfseries,
    title={\rl{🔴 موضوع اختلافی: اسلام‌گرایی چپ است یا راست؟}}
]
\begin{multicols}{2}
\textbf{اسلام‌گرایی «راست» است چون:}
\begin{itemize}
    \item محافظه‌کاری اجتماعی شدید
    \item سلسله‌مراتب جنسیتی
    \item ضد سکولاریسم
    \item سنت‌گرایی
    \item ضد چپ تاریخی
\end{itemize}

\columnbreak

\textbf{اسلام‌گرایی «چپ» است چون:}
\begin{itemize}
    \item ضد امپریالیسم
    \item عدالت اجتماعی در اسلام
    \item ضد سرمایه‌داری غربی
    \item حمایت از فقرا (زکات، صدقه)
    \item ضد نخبگان سکولار
\end{itemize}
\end{multicols}

\vspace{3mm}
\textbf{نتیجه:} اسلام‌گرایی در چارچوب چپ-راست غربی نمی‌گنجد. این جنبش ترکیبی از عناصر مختلف است و شاید نیازمند چارچوب تحلیلی متفاوتی باشد.
\end{tcolorbox}

\begin{multicols}{2}
\textbf{اخوان‌المسلمین:}
تأسیس ۱۹۲۸ در مصر توسط حسن‌البنا. شعار: «اسلام راه‌حل است.» این جنبش هم خدمات اجتماعی ارائه می‌داد و هم خواستار حکومت اسلامی بود.

\textbf{تنوع اسلام‌گرایی:}
\begin{itemize}
    \item \textbf{اسلام‌گرایی اصلاح‌طلب:} مشارکت در دموکراسی (النهضه تونس، AKP ترکیه)
    \item \textbf{اسلام‌گرایی انقلابی:} سرنگونی نظام (جهادی‌ها)
    \item \textbf{سلفی‌گری:} بازگشت به اسلام اولیه (وهابیت)
\end{itemize}

این تنوع نشان می‌دهد که اسلام‌گرایی یک پدیدهٔ واحد نیست.
\end{multicols}

\subsection{ترکیه: از کمالیسم تا اردوغانیسم}

\begin{center}
\begin{tikzpicture}[
    scale=0.85,
    every node/.style={font=\small}
]
% Timeline base
\draw[line width=2pt, color=darkgray] (0,0) -- (16,0);

% Period markers
\foreach \x/\year in {0/1923, 4/1950, 8/1980, 12/2002, 16/2023} {
    \draw[line width=1.5pt, color=darkgray] (\x,-0.2) -- (\x,0.2);
    \node[below] at (\x,-0.3) {\year};
}

% Events
\node[above, align=center, fill=rightblue!20, rounded corners, inner sep=3pt] at (0,0.5) {
    \begin{tabular}{c}
    {\footnotesize\rl{جمهوری ترکیه}}\\
    {\footnotesize\rl{آتاتورک}}
    \end{tabular}
};

\node[above, align=center, fill=goldaccent!20, rounded corners, inner sep=3pt] at (4,1.5) {
    \begin{tabular}{c}
    {\footnotesize\rl{چندحزبی}}\\
    {\footnotesize\rl{حزب دموکرات}}
    \end{tabular}
};

\node[above, align=center, fill=leftred!20, rounded corners, inner sep=3pt] at (8,0.5) {
    \begin{tabular}{c}
    {\footnotesize\rl{کودتا}}\\
    {\footnotesize\rl{نظامیان}}
    \end{tabular}
};

\node[above, align=center, fill=centergreen!20, rounded corners, inner sep=3pt] at (12,1.5) {
    \begin{tabular}{c}
    {\footnotesize\rl{AKP}}\\
    {\footnotesize\rl{اردوغان}}
    \end{tabular}
};

\node[above, align=center, fill=violet!20, rounded corners, inner sep=3pt] at (16,0.5) {
    \begin{tabular}{c}
    {\footnotesize\rl{اقتدارگرایی}}\\
    {\footnotesize\rl{رقابتی}}
    \end{tabular}
};

% Title
\node[above, font=\bfseries\large] at (8,3) {\rl{تایم‌لاین: تحولات سیاسی ترکیه}};

\end{tikzpicture}
\end{center}

\begin{multicols}{2}
\textbf{کمالیسم:}
مصطفی کمال آتاتورک (۱۸۸۱-۱۹۳۸) ترکیهٔ مدرن را بنیان گذاشت:
\begin{itemize}
    \item سکولاریسم رادیکال
    \item ناسیونالیسم ترک
    \item غربی‌سازی فرهنگی
    \item دولت‌گرایی اقتصادی
\end{itemize}

کمالیسم چپ بود یا راست؟ سکولار و مدرنیست (چپ)، اما ناسیونالیست و دولت‌گرای اقتدارگرا.

\textbf{اردوغانیسم:}
رجب طیب اردوغان از ۲۰۰۳ بر ترکیه حاکم است:
\begin{itemize}
    \item اسلام‌گرایی نرم در ابتدا
    \item اقتصاد بازار
    \item ناسیونالیسم ترک
    \item اقتدارگرایی فزاینده
\end{itemize}

AKP (حزب عدالت و توسعه) خود را «دموکرات محافظه‌کار» می‌نامد، اما منتقدان آن را «اقتدارگرایی رقابتی» می‌دانند.
\end{multicols}

%═══════════════════════════════════════════════════════════════
\section{ایران}
%═══════════════════════════════════════════════════════════════

\subsection{از مشروطه تا انقلاب}

\begin{multicols}{2}
\textbf{انقلاب مشروطه (۱۲۸۵):}
اولین انقلاب خاورمیانه برای استقرار پارلمان و قانون اساسی. ترکیبی از:
\begin{itemize}
    \item روشنفکران متجدد
    \item بازاریان سنتی
    \item برخی روحانیون
\end{itemize}

\textbf{چپ ایرانی:}
حزب توده (تأسیس ۱۳۲۰) مهم‌ترین حزب چپ ایران بود. این حزب وابسته به شوروی بود و در کودتای ۱۳۳۲ علیه مصدق نقشی منفعل داشت.

جبههٔ ملی و نهضت آزادی نمایندهٔ ناسیونالیسم لیبرال بودند. \textbf{محمد مصدق} (۱۸۸۲-۱۹۶۷) نماد این جریان بود که در ۱۳۳۲ با کودتای آمریکایی-بریتانیایی سرنگون شد.

سازمان‌های چریکی (فدائیان خلق، مجاهدین خلق) در دههٔ ۵۰ فعال شدند و ترکیبی از مارکسیسم و گاه اسلام ارائه کردند.
\end{multicols}

\subsection{انقلاب اسلامی: فراتر از طیف؟}

\begin{tcolorbox}[
    enhanced,
    colback=goldaccent!10,
    colframe=goldaccent,
    boxrule=1.5pt,
    arc=5pt,
    fonttitle=\bfseries,
    title={\rl{🟡 تحلیل: انقلاب ۱۳۵۷ در طیف سیاسی}}
]
انقلاب ۱۳۵۷ را چگونه می‌توان در طیف چپ-راست جای داد؟

\textbf{عناصر «چپ»:}
\begin{itemize}
    \item ضد امپریالیسم (ضد آمریکا)
    \item شعارهای عدالت اجتماعی («مستضعفین»)
    \item ملی‌سازی صنایع
    \item اصلاحات ارضی (اولیه)
    \item ضد سلطنت و اشرافیت
\end{itemize}

\textbf{عناصر «راست»:}
\begin{itemize}
    \item محافظه‌کاری اجتماعی شدید
    \item تئوکراسی (ولایت فقیه)
    \item سرکوب چپ‌ها و لیبرال‌ها
    \item ضد سکولاریسم
    \item مالکیت خصوصی محترم (با محدودیت)
\end{itemize}

\textbf{نتیجه:} انقلاب اسلامی در چارچوب غربی نمی‌گنجد. شاید بهتر است آن را «انقلاب ضد مدرنیتهٔ غربی» بنامیم تا چپ یا راست.
\end{tcolorbox}

\subsection{جناح‌بندی در جمهوری اسلامی}

\begin{center}
\begin{tikzpicture}[scale=0.9, every node/.style={font=\small}]

% Axes
\draw[line width=2pt, ->] (0,0) -- (12,0) node[right] {\rl{اقتصاد}};
\draw[line width=2pt, ->] (0,0) -- (0,7) node[above] {\rl{سیاست/فرهنگ}};

% Labels on axes
\node[below] at (1,0) {\footnotesize\rl{دولتی}};
\node[below] at (11,0) {\footnotesize\rl{بازار}};
\node[left] at (0,1) {\footnotesize\rl{محافظه‌کار}};
\node[left] at (0,6) {\footnotesize\rl{اصلاح‌طلب}};

% Political factions
\node[draw=leftred, fill=leftred!20, rounded corners, minimum width=2.5cm, minimum height=1cm] at (3,2) {\rl{اصولگرای سنتی}};

\node[draw=rightblue, fill=rightblue!20, rounded corners, minimum width=2.5cm, minimum height=1cm] at (9,2) {\rl{راست سنتی}};

\node[draw=centergreen, fill=centergreen!20, rounded corners, minimum width=2.5cm, minimum height=1cm] at (5,5) {\rl{اصلاح‌طلبان}};

\node[draw=violet, fill=violet!20, rounded corners, minimum width=2.5cm, minimum height=1cm] at (3,3.5) {\rl{چپ اسلامی}};

\node[draw=goldaccent, fill=goldaccent!20, rounded corners, minimum width=2.5cm, minimum height=1cm] at (7,3) {\rl{اعتدالیون}};

% Title
\node[above, font=\bfseries\large] at (6,7.5) {\rl{نقشهٔ جناح‌های سیاسی در جمهوری اسلامی}};

\end{tikzpicture}
\end{center}

\begin{multicols}{2}
\textbf{چپ اسلامی (خط امام):}
در دههٔ ۶۰ غالب. دولتی‌گرا در اقتصاد، رادیکال در سیاست خارجی. میرحسین موسوی نخست‌وزیر این جریان بود.

\textbf{راست سنتی:}
طرفدار بازار و بازاریان. محافظه‌کار در فرهنگ. جامعهٔ روحانیت مبارز.

\textbf{اصلاح‌طلبان:}
از دههٔ ۷۰. خواستار گشایش سیاسی و فرهنگی. محمد خاتمی نماد این جریان.

\textbf{اصولگرایان:}
پس از ۱۳۸۴. محمود احمدی‌نژاد و سپس ابراهیم رئیسی. ترکیبی از پوپولیسم اقتصادی و محافظه‌کاری فرهنگی.

این جناح‌بندی با طیف غربی متفاوت است. همهٔ جناح‌ها در چارچوب نظام هستند و اختلافات در درون این چارچوب است.
\end{multicols}

%═══════════════════════════════════════════════════════════════
\section{آسیای جنوبی: هند}
%═══════════════════════════════════════════════════════════════

\begin{multicols}{2}
\textbf{هند} بزرگ‌ترین دموکراسی جهان است و طیف سیاسی پیچیده‌ای دارد.

\textbf{حزب کنگره:}
حزب استقلال هند. از نهرو تا راهول گاندی. تاریخاً:
\begin{itemize}
    \item سکولاریسم (جدایی دین و دولت)
    \item سوسیالیسم نهرویی (دولتی‌گرایی)
    \item تکثرگرایی (همهٔ ادیان و قومیت‌ها)
\end{itemize}

\textbf{BJP (حزب مردم هند):}
از ۱۹۸۰. حزب حاکم فعلی با نارندرا مودی. ایدئولوژی:
\begin{itemize}
    \item هیندوتوا (ناسیونالیسم هندو)
    \item محافظه‌کاری فرهنگی
    \item اقتصاد بازار (با ملی‌گرایی)
\end{itemize}

\textbf{هیندوتوا} ایدهٔ «هند برای هندوها» است. منتقدان آن را «فاشیسم هندو» می‌نامند؛ مدافعان می‌گویند این صرفاً هویت فرهنگی است.

\textbf{احزاب چپ:}
حزب کمونیست هند (مارکسیست) در برخی ایالات (کرالا، بنگال غربی) قدرت داشته است.

\textbf{احزاب منطقه‌ای و کاستی:}
بخش بزرگی از سیاست هند بر اساس کاست، منطقه، و زبان است—نه چپ و راست.
\end{multicols}

%═══════════════════════════════════════════════════════════════
\section{آسیای شرقی}
%═══════════════════════════════════════════════════════════════

\subsection{چین: از مائو تا شی}

\begin{multicols}{2}
\textbf{مائو تسه‌تونگ} (۱۸۹۳-۱۹۷۶) چین کمونیست را بنیان گذاشت. دوران مائو:
\begin{itemize}
    \item انقلاب ۱۹۴۹
    \item جهش بزرگ (فاجعهٔ اقتصادی)
    \item انقلاب فرهنگی (سرکوب و هرج‌ومرج)
\end{itemize}

\textbf{دنگ شیائوپینگ} (۱۹۰۴-۱۹۹۷) از ۱۹۷۸ اصلاحات را آغاز کرد:
\begin{itemize}
    \item اقتصاد بازار سوسیالیستی
    \item گشایش به سرمایهٔ خارجی
    \item مناطق ویژهٔ اقتصادی
    \item سرکوب تیانانمن (۱۹۸۹)
\end{itemize}

\textbf{شی جین‌پینگ} (از ۲۰۱۲):
\begin{itemize}
    \item تمرکز قدرت
    \item ناسیونالیسم چینی
    \item سرکوب اویغورها و هنگ‌کنگ
    \item رقابت با آمریکا
\end{itemize}

آیا چین امروز «چپ» است؟ حزب کمونیست حاکم است، اما اقتصاد سرمایه‌دارانه است. شاید «اقتدارگرایی سرمایه‌دارانه» دقیق‌تر باشد.
\end{multicols}

\subsection{ژاپن و کرهٔ جنوبی}

\begin{multicols}{2}
\textbf{ژاپن:}
\textbf{LDP} (حزب لیبرال دموکرات) از ۱۹۵۵ تقریباً همیشه حاکم بوده. این حزب:
\begin{itemize}
    \item محافظه‌کار اجتماعی
    \item دولتی‌گرا در توسعهٔ اقتصادی
    \item اتحاد با آمریکا
    \item ملی‌گرایی نسبتاً خفیف
\end{itemize}

اپوزیسیون پراکنده است: حزب دموکراتیک (میانه)، کمونیست (چپ)، و دیگران.

\textbf{کرهٔ جنوبی:}
شکاف اصلی:
\begin{itemize}
    \item \textbf{محافظه‌کاران:} ضد کرهٔ شمالی، اتحاد با آمریکا، توسعهٔ اقتصادی
    \item \textbf{پیشرو‌ها:} گفتگو با شمال، عدالت اجتماعی، حقوق کارگری
\end{itemize}

تاریخ دیکتاتوری نظامی (۱۹۶۱-۱۹۸۷) شکاف عمیقی در جامعه ایجاد کرده.
\end{multicols}

%═══════════════════════════════════════════════════════════════
\section{آمریکای لاتین}
%═══════════════════════════════════════════════════════════════

\subsection{پوپولیسم لاتین}

\begin{multicols}{2}
آمریکای لاتین زادگاه شکلی خاص از پوپولیسم است که نه کاملاً چپ است نه راست.

\textbf{پرونیسم:}
\textbf{خوان پرون} (۱۸۹۵-۱۹۷۴) سه بار رئیس‌جمهور آرژانتین شد. ایدئولوژی‌اش «خوستیسیالیسم» نام داشت:
\begin{itemize}
    \item حمایت از کارگران
    \item ناسیونالیسم اقتصادی
    \item ضد امپریالیسم
    \item «راه سوم» بین سرمایه‌داری و کمونیسم
\end{itemize}

پرونیسم هم جناح چپ دارد و هم راست. این انعطاف‌پذیری ایدئولوژیک، ویژگی پوپولیسم لاتین است.
\end{multicols}

\subsection{موج صورتی و چاوز}

\begin{multicols}{2}
در دهه‌های ۱۹۹۰ و ۲۰۰۰، موجی از دولت‌های چپ در آمریکای لاتین به قدرت رسیدند («موج صورتی»):

\textbf{هوگو چاوز} (۱۹۵۴-۲۰۱۳) در ونزوئلا:
\begin{itemize}
    \item «سوسیالیسم قرن بیست‌ویکم»
    \item ملی‌سازی نفت
    \item برنامه‌های اجتماعی گسترده
    \item ضد آمریکا
    \item اقتدارگرایی فزاینده
\end{itemize}

\textbf{دیگر دولت‌های چپ:}
\begin{itemize}
    \item لولا در برزیل (میانه‌روتر)
    \item اوو مورالس در بولیوی (اولین رئیس‌جمهور بومی)
    \item رافائل کوره‌آ در اکوادور
\end{itemize}

\textbf{تفاوت‌ها:}
\begin{itemize}
    \item چپ میانه (لولا): پذیرش بازار با بازتوزیع
    \item چپ رادیکال (چاوز): ضد سرمایه‌داری
\end{itemize}
\end{multicols}

\subsection{راست جدید: بولسونارو}

\begin{multicols}{2}
در دههٔ ۲۰۱۰، موجی از راست در آمریکای لاتین ظهور کرد.

\textbf{ژائیر بولسونارو} (رئیس‌جمهور برزیل ۲۰۱۹-۲۰۲۲):
\begin{itemize}
    \item «ترامپ برزیل»
    \item محافظه‌کاری اجتماعی شدید
    \item ضد محیط زیست
    \item نظامی‌گرایی
    \item ضد چپ و لولا
\end{itemize}

\textbf{خاویر میلی} (آرژانتین، از ۲۰۲۳):
\begin{itemize}
    \item لیبرتارین افراطی
    \item حذف بانک مرکزی
    \item دلاری‌سازی اقتصاد
    \item ضد دولت رادیکال
\end{itemize}

این نشان می‌دهد که طیف سیاسی آمریکای لاتین همچنان در نوسان است.
\end{multicols}

%═══════════════════════════════════════════════════════════════
\section*{خلاصهٔ فصل}
%═══════════════════════════════════════════════════════════════

\begin{tcolorbox}[
    enhanced,
    colback=lightbg,
    colframe=darkgray,
    boxrule=2pt,
    arc=8pt,
    fonttitle=\bfseries\large,
    title={\rl{📝 خلاصهٔ فصل شانزدهم}}
]
\begin{itemize}
    \item مفاهیم \textbf{چپ و راست} تا حدی در جهان غیرغربی کاربرد دارند، اما باید با شکاف‌های محلی (دین، قوم، استعمار) ترکیب شوند.
    
    \item در \textbf{خاورمیانه}، ناسیونالیسم عربی (ناصریسم، بعثیسم) ترکیبی از چپ اقتصادی و اقتدارگرایی بود. \textbf{اسلام‌گرایی} در چارچوب چپ-راست نمی‌گنجد.
    
    \item در \textbf{ایران}، انقلاب ۱۳۵۷ ترکیبی از عناصر چپ (ضد امپریالیسم) و راست (تئوکراسی) بود. جناح‌بندی داخلی متفاوت از طیف غربی است.
    
    \item در \textbf{هند}، شکاف اصلی بین سکولاریسم (کنگره) و ناسیونالیسم هندو (BJP) است. کاست و منطقه هم نقش مهمی دارند.
    
    \item \textbf{چین} امروز نه سوسیالیست است نه سرمایه‌دار کلاسیک؛ «اقتدارگرایی سرمایه‌دارانه» دقیق‌تر است.
    
    \item \textbf{آمریکای لاتین} تاریخ طولانی پوپولیسم دارد. «موج صورتی» چپ جای خود را به راست جدید (بولسونارو، میلی) داده است.
    
    \item در همه‌جا، شکاف‌های محلی (دین، قوم، استعمار، توسعه) بر شکاف چپ-راست سایه می‌اندازند.
\end{itemize}
\end{tcolorbox}

%═══════════════════════════════════════════════════════════════
\section*{پرسش‌های تأملی}
%═══════════════════════════════════════════════════════════════

\begin{tcolorbox}[
    enhanced,
    colback=centergreen!8,
    colframe=centergreen,
    boxrule=1.5pt,
    arc=5pt,
    fonttitle=\bfseries,
    title={\rl{❓ پرسش‌هایی برای تأمل و بحث}}
]
\begin{enumerate}
    \item آیا تحمیل چارچوب چپ-راست بر جهان غیرغربی نوعی «امپریالیسم فکری» است؟ یا این مفاهیم واقعاً جهانشمول‌اند؟
    
    \item اسلام‌گرایی را چگونه در طیف سیاسی جای دهیم؟ آیا نیازمند چارچوب تحلیلی جدیدی هستیم؟
    
    \item چرا بسیاری از جنبش‌های «چپ» در جهان سوم به اقتدارگرایی انجامیدند؟ آیا این ذاتی چپ است یا شرایط خاص؟
    
    \item آیا «ارزش‌های آسیایی» (که لی کوان یو و دیگران مطرح کردند) واقعی‌اند، یا توجیهی برای اقتدارگرایی؟
    
    \item چرا پوپولیسم در آمریکای لاتین این‌قدر قوی است؟ آیا این به نابرابری شدید مربوط است؟
    
    \item جناح‌بندی در جمهوری اسلامی ایران چه شباهت‌ها و تفاوت‌هایی با طیف سیاسی غربی دارد؟
\end{enumerate}
\end{tcolorbox}

\vspace{10mm}
\begin{center}
\rule{0.6\textwidth}{1pt}

\vspace{3mm}
\textit{«وقتی مفاهیم غربی را به شرق می‌بریم، شرق آنها را می‌بلعد و چیز دیگری پس می‌دهد.»}

\vspace{2mm}
--- ادوارد سعید

\vspace{3mm}
\rule{0.6\textwidth}{1pt}
\end{center}
